\chapter{Deep Dive into Utilized AWS Services}
\label{chap:aws_services}

This chapter provides a detailed, 13-aspect technical breakdown of each AWS service utilized in the College Campus Management System architecture, followed by an explicit evaluation of why Amazon EC2 was not selected.

\section{Amazon Simple Storage Service (Amazon S3)}
\begin{enumerate}
    \item \textbf{Definition}: Object storage service offering industry-leading scalability, data availability, security, and performance.
    \item \textbf{Existence Rationale}: Provides durable, cheap storage for static assets (React bundle HTML/JS/CSS) and dynamic user documents (assignment PDFs, syllabi).
    \item \textbf{Internal Working}: Stores data as objects within buckets. Each object consists of data bytes, a key, and metadata, replicated across at least three physical AZs.
    \item \textbf{Problem Solved}: Eliminates disk space constraints and disk failures associated with local web server storage.
    \item \textbf{Project Usage}: Hosts the frontend Single Page Application via Static Website Hosting and stores uploaded student assignment files in \texttt{/uploads}.
    \item \textbf{Configuration Used}: Bucket name \texttt{cloudcampus-frontend-static}, public read access via S3 Bucket Policy for static assets, CORS rule allowing \texttt{GET, POST, PUT} from application domain.
    \item \textbf{Data Flow}: Receives compiled static assets during deployment and receives assignment binary streams uploaded by students via Lambda proxy.
    \item \textbf{Inter-service Connectivity}: Integrated with API Gateway and Lambda via AWS SDK (\texttt{@aws-sdk/client-s3}).
    \item \textbf{Security}: Server-side encryption with Amazon S3 managed keys (SSE-S3, AES-256), Block Public Access enforced on raw storage buckets.
    \item \textbf{Cost Model}: Charges \$0.023/GB/month for standard storage and \$0.0004 per 1,000 PUT/POST requests.
    \item \textbf{Observations}: Achieved sub-50ms latency for static web asset retrieval.
    \item \textbf{Advantages}: Infinite scaling, 99.999999999\% durability.
    \item \textbf{Limitations}: High GET request costs if un-cached by a CDN.
\end{enumerate}

\section{Amazon Relational Database Service (Amazon RDS for PostgreSQL)}
\begin{enumerate}
    \item \textbf{Definition}: Fully managed relational database service supporting PostgreSQL 15.3.
    \item \textbf{Existence Rationale}: Handles complex relational queries, foreign key constraints, ACID transactions, and structured academic data.
    \item \textbf{Internal Working}: Runs PostgreSQL engine inside a dedicated VPC private subnet with EBS volume storage and automated WAL archiving.
    \item \textbf{Problem Solved}: Eliminates manual database patching, replication setup, and storage management.
    \item \textbf{Project Usage}: Stores system records for Users, Students, Faculty, Departments, Courses, Enrollments, Attendance, Assignments, Results, Events, Announcements, Notifications, and Audit Logs.
    \item \textbf{Configuration Used}: PostgreSQL 15.3, \texttt{db.t4g.micro} instance class, 20 GB Storage, VPC Private Subnet placement, Security Group restricting ingress to Lambda security group on port 5432/5433.
    \item \textbf{Data Flow}: Receives SQL queries executed via Prisma Client originating from AWS Lambda backend functions.
    \item \textbf{Inter-service Connectivity}: Connected to AWS Lambda via VPC Security Group ingress rules.
    \item \textbf{Security}: Encrypted at rest using AWS KMS, TLS-encrypted connections enforced, restricted to private subnets.
    \item \textbf{Cost Model}: Hourly instance rate (\$0.017/hr for \texttt{db.t4g.micro}) + \$0.115/GB/month allocated storage.
    \item \textbf{Observations}: Query response times averaged 4–12ms across index-optimized queries.
    \item \textbf{Advantages}: Automated daily backups, zero-downtime storage scaling.
    \item \textbf{Limitations}: Connection pool limit must be managed carefully when paired with serverless Lambda.
\end{enumerate}

\section{AWS Lambda}
\begin{enumerate}
    \item \textbf{Definition}: Serverless, event-driven compute service executing code in response to events without server management.
    \item \textbf{Existence Rationale}: Runs Express REST API microservices dynamically, scaling instantly from zero to high request concurrency.
    \item \textbf{Internal Working}: Provisions ephemeral MicroVM containers (Firecracker) executing the Node.js 20.x runtime on demand.
    \item \textbf{Problem Solved}: Removes server maintenance and eliminates costs for idle server execution.
    \item \textbf{Project Usage}: Executes API logic for authentication, course enrollment, attendance logging, grading, and report generation.
    \item \textbf{Configuration Used}: Node.js 20.x runtime, 512 MB memory allocation, 15-second execution timeout, attached to VPC private subnets.
    \item \textbf{Data Flow}: Ingests JSON payload events from API Gateway, executes business logic, queries RDS PostgreSQL, and returns JSON responses.
    \item \textbf{Inter-service Connectivity}: Triggered by API Gateway; reads/writes to RDS, S3, SNS, and CloudWatch.
    \item \textbf{Security}: Execution role bounded by least-privilege IAM policies.
    \item \textbf{Cost Model}: \$0.20 per 1,000,000 requests + \$0.0000166667 per GB-second of execution.
    \item \textbf{Observations}: Cold starts averaged 240ms; warm execution duration averaged 18–35ms.
    \item \textbf{Advantages}: Automatic scaling, pay-per-use billing down to the millisecond.
    \item \textbf{Limitations}: Execution time limit capped at 15 minutes.
\end{enumerate}

\section{Amazon API Gateway}
\begin{enumerate}
    \item \textbf{Definition}: Fully managed service for creating, publishing, maintaining, monitoring, and securing REST and HTTP APIs at scale.
    \item \textbf{Existence Rationale}: Exposes secure public HTTP endpoints for the React SPA and proxies incoming traffic to AWS Lambda.
    \item \textbf{Internal Working}: Receives HTTP requests at regional edge locations, performs routing, validates headers, and invokes target Lambda handlers.
    \item \textbf{Problem Solved}: Handles web request routing, SSL termination, and rate limiting out of the box.
    \item \textbf{Project Usage}: Manages 44 REST API routes under \texttt{/api/*} mapped to backend Lambda handlers.
    \item \textbf{Configuration Used}: REST API, HTTP proxy integration with Lambda, CORS enabled, throttling limit of 200 requests per 15-minute window.
    \item \textbf{Data Flow}: Receives HTTPS client requests, proxies request context to Lambda, and returns HTTP responses to frontend.
    \item \textbf{Inter-service Connectivity}: Integrates directly with Lambda and Cognito Authorizers.
    \item \textbf{Security}: TLS 1.3 encryption, CORS domain restriction, Cognito JWT authorization.
    \item \textbf{Cost Model}: \$3.50 per million API calls.
    \item \textbf{Observations}: Zero request drops observed under concurrent client testing.
    \item \textbf{Advantages}: Seamless integration with Cognito and Lambda; native rate limiting.
    \item \textbf{Limitations}: 29-second payload timeout.
\end{enumerate}

\section{Amazon Cognito}
\begin{enumerate}
    \item \textbf{Definition}: Managed identity service providing user sign-up, sign-in, and access control.
    \item \textbf{Existence Rationale}: Manages user credentials securely and issues standard JSON Web Tokens (JWT) for API authentication.
    \item \textbf{Internal Working}: Maintains user pools, hashes passwords using bcrypt/PBKDF2, and issues signed RSA-256 JWT tokens.
    \item \textbf{Problem Solved}: Eliminates custom password hashing and security risks associated with storing sensitive credentials.
    \item \textbf{Project Usage}: Authenticates Student, Faculty, and Admin users and issues bearer tokens.
    \item \textbf{Configuration Used}: Cognito User Pool \texttt{cloudcampus-user-pool}, custom attributes for Role (\texttt{STUDENT}, \texttt{FACULTY}, \texttt{ADMIN}).
    \item \textbf{Data Flow}: Accepts user credentials during login, validates against pool, returns JWT access and ID tokens to frontend.
    \item \textbf{Inter-service Connectivity}: Authenticates requests before API Gateway forwards them to Lambda.
    \item \textbf{Security}: MFA support, password strength enforcement, signed RSA tokens.
    \item \textbf{Cost Model}: Free tier includes 50,000 Monthly Active Users (MAUs).
    \item \textbf{Observations}: Token validation completed in under 5ms at API Gateway.
    \item \textbf{Advantages}: Standards-compliant OAuth 2.0 / OIDC implementation.
    \item \textbf{Limitations}: Custom UI authorization flows require careful token parsing.
\end{enumerate}

\section{Amazon Simple Notification Service (Amazon SNS)}
\begin{enumerate}
    \item \textbf{Definition}: Managed pub/sub messaging service for asynchronous message delivery.
    \item \textbf{Existence Rationale}: Dispatches instant academic alerts and system announcements to students and faculty.
    \item \textbf{Internal Working}: Receives messages on topics and pushes them to subscribed endpoints (Email, SMS, Lambda).
    \item \textbf{Problem Solved}: Decouples notification sending from core API request execution threads.
    \item \textbf{Project Usage}: Sends urgent exam alerts, assignment grading notifications, and campus announcements.
    \item \textbf{Configuration Used}: Standard SNS Topic \texttt{cloudcampus-notifications}, Email subscription protocol.
    \item \textbf{Data Flow}: Lambda publishes notification events to SNS topic; SNS broadcasts messages to subscribers.
    \item \textbf{Inter-service Connectivity}: Published to by Lambda backend functions.
    \item \textbf{Security}: Topic access controlled via IAM policies and SNS access control policies.
    \item \textbf{Cost Model}: \$0.50 per 1,000,000 SNS requests.
    \item \textbf{Observations}: Email notifications delivered within 1.2 to 2.5 seconds of trigger.
    \item \textbf{Advantages}: Reliable asynchronous push delivery.
    \item \textbf{Limitations}: Email delivery depends on recipient mail server filters.
\end{enumerate}

\section{Amazon CloudWatch}
\begin{enumerate}
    \item \textbf{Definition}: Monitoring and observability service providing logs, metrics, and event alerts.
    \item \textbf{Existence Rationale}: Captures application logs, tracks Lambda latency, and monitors database health.
    \item \textbf{Internal Working}: Collects log streams from Lambda and RDS, aggregates metrics, and evaluates alarm rules.
    \item \textbf{Problem Solved}: Centralizes observability across distributed serverless components.
    \item \textbf{Project Usage}: Stores execution log group \texttt{/aws/lambda/cloudcampus-backend-api} and monitors API Gateway latency.
    \item \textbf{Configuration Used}: 14-day log retention policy, custom metrics for 5xx errors.
    \item \textbf{Data Flow}: Automatically ingests \texttt{console.log} and \texttt{console.error} streams from Lambda runtime containers.
    \item \textbf{Inter-service Connectivity}: Receives telemetry from Lambda, API Gateway, RDS, and VPC.
    \item \textbf{Security}: Encryption at rest using KMS; access restricted via IAM permissions.
    \item \textbf{Cost Model}: \$0.50 per GB log ingestion; free tier includes 5 GB ingestion.
    \item \textbf{Observations}: Enabled rapid root-cause diagnosis of database connection timeouts.
    \item \textbf{Advantages}: Native integration across every AWS resource.
    \item \textbf{Limitations}: Log ingestion costs can escalate if verbose debugging is left enabled.
\end{enumerate}

\section{AWS Identity and Access Management (AWS IAM)}
\begin{enumerate}
    \item \textbf{Definition}: Fine-grained access control service defining who can perform actions on AWS resources.
    \item \textbf{Existence Rationale}: Enforces least-privilege security permissions across all cloud services.
    \item \textbf{Internal Working}: Evaluates JSON policy documents attached to users, roles, and resources for authorization.
    \item \textbf{Problem Solved}: Prevents unauthorized service access and eliminates hardcoded credentials.
    \item \textbf{Project Usage}: Grants Lambda permissions to read/write S3 objects, publish to SNS topics, and write logs to CloudWatch.
    \item \textbf{Configuration Used}: Execution role \texttt{CloudCampusLambdaExecutionRole} with policies \texttt{AWSLambdaVPCAccessExecutionRole} and custom inline policy for S3/SNS access.
    \item \textbf{Data Flow}: Evaluates temporary security credentials issued via AWS Security Token Service (STS).
    \item \textbf{Inter-service Connectivity}: Underpins authorization across every AWS service call.
    \item \textbf{Security}: Enforces strict policy syntax with explicit default deny.
    \item \textbf{Cost Model}: Free service provided with AWS accounts.
    \item \textbf{Observations}: Completely eliminated credential leaks from application code.
    \item \textbf{Advantages}: Fine-grained resource-level permission granularity.
    \item \textbf{Limitations}: Policy syntax complexity requires careful auditing.
\end{enumerate}

\section{Amazon Virtual Private Cloud (Amazon VPC)}
\begin{enumerate}
    \item \textbf{Definition}: Isolated virtual network dedicated to an AWS account.
    \item \textbf{Existence Rationale}: Isolates the PostgreSQL database and backend Lambda compute from public internet exposure.
    \item \textbf{Internal Working}: Provisions private IP subnets, route tables, network ACLs, and security groups on AWS hardware.
    \item \textbf{Problem Solved}: Prevents direct internet scanning and brute-force attacks against database ports.
    \item \textbf{Project Usage}: Houses RDS database instance in private subnets across two AZs.
    \item \textbf{Configuration Used}: VPC CIDR \texttt{10.0.0.0/16}, 2 Public Subnets, 2 Private Subnets, Security Group \texttt{sg-cloudcampus-db} allowing TCP 5432 from Lambda.
    \item \textbf{Data Flow}: Filters network traffic at security group and network ACL boundaries.
    \item \textbf{Inter-service Connectivity}: Encloses Lambda VPC ENIs and RDS instances.
    \item \textbf{Security}: Network ACLs, stateful security groups, zero public IP assignment to private subnets.
    \item \textbf{Cost Model}: Standard VPC features are free; NAT Gateways incur hourly charges if provisioned.
    \item \textbf{Observations}: Ensured zero unauthorized port scans reached the database layer.
    \item \textbf{Advantages}: Complete network isolation and software-defined firewalling.
    \item \textbf{Limitations}: Cold start delay incurred when Lambda creates Elastic Network Interfaces (ENIs) inside VPC.
\end{enumerate}

\section{Why Amazon EC2 Was Not Selected}
Amazon Elastic Compute Cloud (Amazon EC2) is AWS's flagship Infrastructure as a Service (IaaS) offering, providing virtual machines (instances) running Linux or Windows operating systems. While EC2 is the conventional choice for hosting traditional web servers, it was intentionally excluded from the final cloud architecture for this project due to specific architectural and financial trade-offs:

\begin{enumerate}
    \item \textbf{Continuous Idle Cost Burden}: An EC2 instance incurs billing for every hour it runs, regardless of whether traffic is being processed. For an academic campus application with low night-time usage, paying for 24/7 idle EC2 compute is financially inefficient compared to AWS Lambda's pay-per-execution model.
    \item \textbf{Administrative Maintenance Overhead}: Deploying an EC2 instance requires administrators to perform operating system security patching, manage package dependencies, configure systemd process supervisors (e.g., PM2), and set up custom firewall rules. Serverless Lambda removes all OS management.
    \item \textbf{Scaling Complexity}: Achieving high availability on EC2 requires configuring Auto Scaling Groups, Application Load Balancers (ALBs), launch templates, and multi-AZ instance deployment. In contrast, AWS Lambda automatically scales concurrency up and down without load balancer configuration.
    \item \textbf{Resource Utilization Efficiency}: EC2 virtual machines assign fixed CPU cores and RAM. Under variable institutional demand, an EC2 instance is either over-provisioned (wasting money) or under-provisioned (risking downtime). Lambda provides exact sub-second resource allocation.
\end{enumerate}

Therefore, excluding EC2 in favor of a serverless architecture (S3 + API Gateway + Lambda + RDS) resulted in a lower cost footprint, superior elasticity, and zero operating system maintenance overhead.
